% BHU PYQ -- Manually transcribed & error-corrected
\documentclass[11pt,a4paper]{article}

\usepackage[a4paper,top=2.5cm,bottom=2.5cm,left=2.8cm,right=2.8cm,headheight=14pt]{geometry}
\usepackage{amsmath,amssymb}
\usepackage[T1]{fontenc}
\usepackage[utf8]{inputenc}
\usepackage{lmodern}
\usepackage{microtype}
\usepackage{fancyhdr}
\usepackage{enumitem}
\usepackage{hyperref}
\usepackage{lastpage}
\usepackage{parskip}

\hypersetup{colorlinks=true,urlcolor=black,linkcolor=black,
  pdftitle={BP-Anc-I: Ancillary Physics-I -- BHU PYQ},pdfauthor={Banaras Hindu University}}

\pagestyle{fancy}\fancyhf{}
\renewcommand{\headrulewidth}{0.4pt}\renewcommand{\footrulewidth}{0.4pt}
\fancyhead[L]{\small   \textit{Banaras Hindu University}}
\fancyhead[R]{\small   \textit{BP-Anc-I: Ancillary Physics-I}}
\fancyfoot[C]{\small Page  \thepage\ of \pageref{LastPage}}


\newlist{parts}{enumerate}{2}
\setlist[parts,1]{label=(\alph*),leftmargin=*,itemsep=5pt,topsep=3pt}
\setlist[parts,2]{label=(\roman*),leftmargin=*,itemsep=3pt,topsep=2pt}

\newcommand{\pts}[1]{\hfill   \textit{[#1]}}


\begin{document}


\begin{center}
  {\large\bfseries BANARAS HINDU UNIVERSITY}\\[2pt]
  {\normalsize B.Sc. (Hons.) Semester II Examination, 2013-14}\\[6pt]
  \rule{0.8\linewidth}{0.5pt}\\[6pt]
  {\large\bfseries Paper No. BP-Anc-I: Ancillary Physics-I}\\[4pt]
  {\normalsize Time: 3 Hours\hfill Maximum Marks: 70}
\end{center}
\rule{\linewidth}{0.4pt}
\smallskip



\noindent   \textit{Instructions: Answer   \textbf{five} questions including Question~1 which is compulsory. Symbols have their usual meanings.}
\bigskip




\noindent   \textbf{Question 1.} Answer any   \textbf{six} of the following: \pts{6 $ \times$ 3 = 18}

\begin{parts}
  \item The electromagnetic radiations in different wavelength/frequency ranges are used for different biological applications and in different biological processes. In each of the following three processes, which region of electromagnetic radiation is used?
    
\begin{parts}
      \item Photosynthesis
      \item For killing bacteria
      \item For killing microbes in a food material
    \end{parts}
  \item Why are some isotopes radioactive and others not? Can you predict which one will be radioactive?
  \item List the cellular sensitivity of radiation in terms of most sensitive, less sensitive, and least sensitive.
  \item Explain some of the important applications of elasticity in biology.
  \item Define and give at least two examples of a perfectly elastic and a perfectly plastic body.
  \item Explain why liquid drops are more nearly spherical as their radius decreases.
  \item What do you know about Laplace's Law?
  \item Why is it better to use hot soapy water to wash clothes?
\end{parts}

\medskip


\noindent   \textbf{Question 2.}

\begin{parts}
  \item Define the coefficient of viscosity and explain, with necessary theory, a method to find the viscosity of a liquid. \pts{5}
  \item Explain the phenomenon of photosynthesis briefly in terms of a simple chemical reaction. \pts{4}
  \item What are the various harmful effects of deep UV radiation on human beings? Describe any two of them in brief. \pts{4}
\end{parts}

\medskip


\noindent   \textbf{Question 3.}

\begin{parts}
  \item Explain the biological effects of radiations. How does ionizing radiation damage DNA? \pts{5}
  \item What are the consequences of high doses of radiation over short periods of time? Explain in brief. \pts{4}
  \item Explain why raindrops falling under gravity do not acquire a very high velocity. \pts{4}
\end{parts}

\medskip


\noindent   \textbf{Question 4.}

\begin{parts}
  \item What do you understand by surface tension and angle of contact? State the condition under which the angle of contact is acute and obtuse. What will be the numerical value of the angle of contact when the liquid is said to wet the solid? What are the factors on which the angle of contact depends upon? What are the typical values of the angle of contact for a pair of:
    
\begin{parts}
      \item clean glass and pure water,
      \item ordinary water and glass,
      \item ordinary water and greased glass surface,
      \item ordinary water and chromium?
    \end{parts} \pts{8}
  \item Explain the inflation mechanism of the alveoli and the role of surfactant in human respiration. We know that it is much more difficult to blow up a balloon for the first time. Why is that? \pts{5}
\end{parts}

\medskip


\noindent   \textbf{Question 5.}

\begin{parts}
  \item What do you understand by excess of pressure inside a liquid drop? Describe in detail Jaeger's experimental method for the determination of surface tension. Explain the procedure, draw a neat sketch of the apparatus, and write precautions. \pts{8}
  \item Explain the behavior of a wire under increasing load with the help of a neat sketch and define yield-point, plastic-flow, and breaking-point. \pts{5}
\end{parts}

\medskip


\noindent   \textbf{Question 6.}

\begin{parts}
  \item A thin, circular wire of diameter $40\, \text{mm}$ and total mass of $0.70\, \text{g}$ is gently pulled vertically from a water surface by a sensitive spring ($k = 0.70\, \text{N/m}$). When the spring is stretched $34\, \text{mm}$ from its equilibrium position in air, the ring is on the verge of being pulled free from the water surface. Find the coefficient of surface tension of water. Neglect the mass of water lifted. \pts{8}
  \item Show the non-linear elasticity behavior of tissue and explain it with the help of a stress-strain diagram. \pts{8}
\end{parts}

\vfill
\rule{\linewidth}{0.4pt}

\begin{center}\small
  \href{https://bhu-exam.vercel.app/}{Click here to visit our website}\\[4pt]\href{https://me-aryan.vercel.app/}{Click here to visit our contributor}\\[4pt]\href{https://quashan-me.vercel.app/}{Click here to visit our contributor}
\end{center}

\end{document}
